\section{观点后验、尾部目标与稳健实施}
\label{ch:lower-portfolio-optimization}

\begin{itemize}
  % Generated from curriculum/manifest.json; do not edit by hand.
\item 直接先修：下册《风险估计：从协方差到风险预算》、上册第 13 章《凸优化、KKT 与对偶》。

\end{itemize}

观点可信多少与持仓承担多少风险是不同输入。先用一个标量后验计算理解信息合并，再用不确定集推导保守收益目标。


\MFQTransition{Black--Litterman：把观点和置信度显式化}

% Generated by tools/build_figures.py; edit figures/figure-specs.json instead.
\MFQFigure{tex/figures/generated/lower-pf-frontier.pdf}{增加约束或计入成本会缩小可达到的收益与风险组合；最优选择也随风险偏好与换手代价变化。}{lower-pf-frontier}


设先验均值 $\pi$ 的不确定性为 $\tau\Sigma$，观点矩阵 $P$、观点收益 $q$、观点误差协方差 $\Omega$ 满足
\[
 P\mu=q+\eta,\qquad \eta\sim N(0,\Omega).
\]
高斯更新给出
\[
 \mu_{\mathrm{BL}}=
 \left[(\tau\Sigma)^{-1}+P^\top\Omega^{-1}P\right]^{-1}
 \left[(\tau\Sigma)^{-1}\pi+P^\top\Omega^{-1}q\right].
\]
后验均值的不确定性为括号中精度矩阵的逆；若预测收益本身也随机，后验预测协方差还要加回 $\Sigma$\cite{black1992global}。

算例的先验均值为 $(0.05,0.04)$，观点是第一资产相对第二资产超额 $0.03$，$\tau=0.05$，观点方差为 $0.0025$。后验均值为 $(0.05375,0.03)$。透明精度公式与独立 Woodbury 更新最大差 $2.8\times10^{-17}$。当 $\Omega$ 变小，观点权重上升；这表示模型更相信观点，不表示观点变得更真实。

\MFQTransition{CVaR 的线性规划形式}

给定 $N$ 个情景收益 $r_s$，损失 $L_s(w)=-r_s^\top w$。Rockafellar--Uryasev 表示为
\[
 \operatorname{CVaR}_\alpha(w)=
 \min_{z}\left[z+\frac{1}{(1-\alpha)N}
 \sum_{s=1}^{N}(L_s(w)-z)_+\right].
\]
引入辅助变量 $u_s\ge0$ 得到线性规划
\[
 \min_{w,z,u}\quad z+\frac{1}{(1-\alpha)N}\sum_su_s,
 \qquad u_s\ge-r_s^\top w-z,\quad
 \mathbf1^\top w=1,\quad 0\le w_i\le\bar w.
\]
这不是把样本 VaR 当作平滑目标，而是同时选择权重和尾部阈值\cite{rockafellar2000optimization}。

事先确定八情景、$\alpha=0.75$、单资产上限 $0.8$ 的 LP 得到 $(0.2,0.8)$，CVaR 为 $0.013$。独立 0.1 网格枚举重算每个候选的最坏两情景均值，得到相同权重和目标。若枚举忘记权重上限，会错误选择第二资产 0.9；这正说明可行域也是 解析参照 的一部分。

\MFQTransition{把估计误差、成本和可交易性放在一起}

稳健再平衡采用
\[
 \max_w\quad
 \widehat\mu^\top w
 -\rho\lVert U w\rVert_2
 -\frac\gamma2w^\top\Sigma w
 -c\lVert w-w^-\rVert_1,
\]
并约束满仓、长仓、最大权重和不可交易资产权重不变。$U$ 编码期望收益不确定性尺度，$\rho$ 是最坏情形惩罚。它与协方差 ridge 都能平滑权重，但经济对象不同：一个约束均值误差，一个稳定风险矩阵。

事先确定结果从 $(0.4,0.6)$ 调整为 $(0.8,0.2)$。双边换手为 $0.8$；十万元资本、每单位换手成本 $0.001$ 的现金成本为
\[
 100000\times0.001\times0.8=80.
\]
比例成本与收益采用同一口径；乘以资本规模后得到相应的现金金额。若机构采用单边换手，必须显式除以二；基点、百分数和小数的混用会造成百倍或万倍错误。

\MFQTransition{均值--方差优化与输入放大}

无约束均值--方差解为
\[
w^*=\frac1\gamma\Sigma^{-1}\mu.
\]
它清楚展示了两个放大器：均值估计误差直接进入权重，协方差小特征值又通过逆矩阵放大。实际系统必须使用权重界、暴露约束、均值收缩和交易成本，而不是把解析解当推荐组合。

若目标是相对基准 $b$ 的主动组合，使用主动权重 $a=w-b$ 和跟踪误差 $a^{\mathsf T}\Sigma a$。行业中性、beta 中性和换手上限都可写成线性/凸约束。不可交易资产的旧仓位要固定，而不是在求解后裁成零。

\MFQTransition{观点的单位与尾部数据}
相对观点 $P=(1,-1)$ 约束的是两资产均值差，绝对观点 $P=(1,0)$ 约束的是第一资产均值。观点误差方差与收益均值采用同一期间。CVaR 的辅助阈值属于损失坐标，情景数少时线性规划仍有解，却可能由极少数观测决定。

\MFQTransition{稳健优化与分布不确定性}

若均值属于椭球不确定集，最坏情形均值会产生额外范数惩罚；协方差不确定可用情景集合或分布鲁棒方法。稳健半径越大，组合越保守。半径应由估计误差或验证协议决定，不能为获得目标仓位反推。

优化报告应包含 shadow price 或活跃约束：它们说明性能被哪个限制约束。若小幅改变上限就导致完全不同组合，说明解位于尖锐边界，应做参数和输入扰动。

\begin{diagnostic}
求解器返回 optimal 只表示数值模型被解出。若预期收益来自未来修订、成本单位错误或场景不代表尾部，最优解仍是错误问题的精确答案。
\end{diagnostic}

优化结果将在组合风险路线综合项目中接受尾部压力、输入扰动和实施限制的共同复核。
\subsection{精度加权与最坏均值}

标量先验 $\mu\sim N(\pi,v)$，观测观点 $q=\mu+\eta$，$\eta\sim N(0,\omega)$ 独立。配方得到后验均值 $m=(\omega\pi+vq)/(v+\omega)$，方差 $v\omega/(v+\omega)$。$\pi=0.04,q=0.08,v=\omega$ 时，$m=0.06$；若观点方差降为 $v/3$，则 $m=0.07$。高置信度提高观点权重，但不检验观点是否正确。

若均值不确定集为 $\mu=\hat\mu+U^Tz$、$\|z\|_2\leq\rho$，则
\[
 \min_\mu\mu^Tw=\hat\mu^Tw+\min_{\|z\|\leq\rho}z^TUw
 =\hat\mu^Tw-\rho\|Uw\|_2.
\]
Cauchy--Schwarz 给下界；当 $Uw\ne0$ 时取 $z=-\rho Uw/\|Uw\|$ 达到下界。这解释了稳健目标中的范数惩罚来自哪里。notebook 显示标量后验随置信度变化，并枚举两资产权重、计算每个权重的尾部损失和成本。
